\section{Introduction}

Modern database management systems are increasingly \emph{self-driving}: they
observe their own workload and act on it without human intervention
\citep{pavlo2017selfdriving}. Physical design---in particular, choosing which
secondary indexes to create and drop---is one of the highest-leverage knobs in
such a system, because indexes can change query latency by orders of magnitude.
Classical index advisors formulate this as a what-if optimization problem over a
fixed, known workload and solve it with greedy search or integer programming
\citep{chaudhuri1997autoadmin,valentin2000db2advisor,zilio2004db2design}. These
advisors are \emph{offline} and \emph{assume a stationary workload}: the workload
is collected, the advisor is run, indexes are materialized, and the loop stops.

Recent work replaces the hand-designed search with learned models, including
learned cost estimators \citep{marcus2019learnedcost}, learned cardinality
estimators \citep{kipf2019learnedcardinalities}, learned index structures
\citep{kraska2018case}, and learned query optimizers such as Bao
\citep{marcus2021bao}. A parallel line applies reinforcement learning (RL) to
physical design, training an agent to propose index actions and rewarding it with
observed latency improvements. All of these systems share the same hidden
assumption as their classical predecessors: \emph{the workload used to train the
model is the workload that will be served at deployment time.}

Production workloads violate this assumption. A transactional system looks very
different at 9~a.m.\ (short, selective, point lookups on the \texttt{orders}
table) than at 2~a.m.\ (long analytical scans over the same table); retail
systems shift with campaigns and seasons; and schemas and access patterns evolve
over months. We call such a temporally evolving workload a \emph{workload drift},
and the periods over which the workload distribution is approximately stationary
\emph{workload phases}. When the workload drifts, a deployed advisor faces a
dilemma:
\begin{enumerate}
  \item \textbf{Offline retraining} from scratch on the new workload is
    expensive, requires collecting and labelling a batch of queries, and
    disrupts the service; it also discards knowledge of previous phases that may
    return.
  \item \textbf{Online fine-tuning} on recent queries is cheap and
    non-disruptive, but naively updating the parameters on the new phase causes
    \emph{catastrophic forgetting} \citep{mccloskey1989catastrophic}: the agent
    rapidly loses the ability to optimize the phases it previously mastered,
    which is exactly the situation when the workload later returns to a previous
    phase.
\end{enumerate}
This is the continual-learning trilemma---\emph{adapt}, \emph{retain}, and
\emph{stay cheap/online}---and existing index advisors address at most two of the
three.

\paragraph{Approach.}
We argue that autonomous index selection should be posed not as a single-task RL
problem but as a \emph{continual reinforcement-learning} problem
\citep{khetarpal2022continual}. We therefore introduce \textbf{CLDB}, an online
index-selection agent that learns an unbounded sequence of non-stationary
workload phases from live query feedback. CLDB makes four design contributions,
all of which are implemented end-to-end:
\begin{itemize}
  \item \textbf{Index-aware off-policy DQN.} The agent is a deep Q-network
    \citep{mnih2015dqn} with a target network. Crucially, the state includes the
    \emph{current index configuration} (number of active indexes, whether the
    current predicate is covered, and remaining budget), making the process a
    genuine Markov decision process rather than a contextual bandit: creating an
    index has a one-time cost and a recurring benefit, and dropping it is a
    decision that depends on the accumulated index state.
  \item \textbf{Drift-triggered EWC.} Instead of consolidating at fixed task
    boundaries (which are unknown online), CLDB computes the elastic weight
    consolidation penalty \citep{kirkpatrick2017ewc} only when a drift detector
    \citep{gama2014drift} fires. We analyze the cost of this event-driven rule
    relative to boundary-based consolidation (Lemma~\ref{lem:consolidation}).
  \item \textbf{Database-aware prioritized replay.} The replay buffer
    prioritizes experiences by their \emph{database impact}---absolute latency
    reward times workload rarity---and refines those priorities with absolute TD
    errors, combining the anti-forgetting effect of experience replay
    \citep{rolnick2019replay} with the sample efficiency of prioritized replay
    \citep{schaul2016prioritized}.
  \item \textbf{Safety-constrained exploration.} A structural action mask forbids
    \texttt{CREATE\_INDEX} without predicate columns and \texttt{DROP\_INDEX}
    without an existing index, so exploration can never propose an unsafe schema
    change.
\end{itemize}

\paragraph{Research questions.}
CLDB is designed to answer the following question and its refinements:
\begin{quote}
\emph{Can a reinforcement-learning index-selection policy continually acquire new
workload phases from live query feedback---while retaining competence on
previously seen phases---without offline retraining or blocking query
execution?}
\end{quote}
Concretely we ask: \textbf{(RQ1)} does the agent retain earlier phases better
than naive online fine-tuning? \textbf{(RQ2)} which mechanism drives retention,
replay or EWC? \textbf{(RQ3)} does database-aware prioritization beat uniform
replay? \textbf{(RQ4)} does action masking help or hurt? \textbf{(RQ5)} can
adaptation happen online at bounded cost?

\paragraph{Contributions.}
\begin{enumerate}
  \item We give the first formulation of index selection as a continual
    reinforcement-learning problem with an index-aware state, and we provide a
    reproducible \emph{cost-model benchmark} of conflicting workload phases that
    induces genuine catastrophic forgetting (Section~\ref{sec:problem},
    Section~\ref{sec:experiments}).
  \item We design and implement a complete agent---drift-triggered EWC,
    database-aware prioritized replay, and safety masking---together with an
    online pipeline that satisfies the non-blocking requirement
    (Section~\ref{sec:method}).
  \item We provide a theoretical analysis: the EWC objective as a Laplace
    posterior approximation (Proposition~\ref{prop:ewc}), the value-network
    Fisher estimator (Proposition~\ref{prop:fisher}), the unbiasedness of the
    prioritized-replay gradient (Proposition~\ref{prop:per}), and the
    consolidation-cost advantage of drift triggering
    (Lemma~\ref{lem:consolidation}) (Section~\ref{sec:theory}).
  \item We evaluate against static, random and heuristic baselines and against
    ablations (no EWC, no replay) using the standard continual-learning metrics
    ACC, BWT, FWT and forgetting \citep{lopezpaz2017gem}, with eight seeds and
    paired significance tests (Section~\ref{sec:results}). The central,
    statistically significant finding is that \emph{rehearsal} (replay) is the
    ingredient that prevents catastrophic forgetting; EWC's incremental benefit is
    smaller and, in this benchmark, not significant---a result we report and
    analyze rather than hide.
\end{enumerate}

\paragraph{Organization.}
Section~\ref{sec:related} reviews related work. Section~\ref{sec:problem} gives
the mathematical formulation. Section~\ref{sec:method} describes CLDB.
Section~\ref{sec:theory} develops the theory. Sections~\ref{sec:experiments}
and~\ref{sec:results} present the setup and results.
Section~\ref{sec:discussion} discusses limitations and threats to validity, and
Section~\ref{sec:conclusion} concludes.
